% !TeX program = xelatex
% 编译：xelatex 附录代码片段.tex
% 依赖：ctexart + listings（MiKTeX / TeX Live 均自带），中文无需额外配置。
% 并入论文：把正文（document 环境内的内容）复制进 A题论文.tex 的附录处，
%           并在论文导言区加入 \usepackage{listings,xcolor} 与本文件的 \lstset 配置。

\documentclass[12pt,a4paper]{ctexart}
\usepackage{amsmath,amssymb}
\usepackage{booktabs}
\usepackage{array}
\usepackage[margin=2.5cm]{geometry}
\usepackage{listings}
\usepackage{xcolor}
\usepackage[hidelinks]{hyperref}

\linespread{1.25}
\setlength{\parskip}{0.5em}
\setlength{\parindent}{2em}

\ctexset{
  section = {format = {\heiti \zihao{3}}},
  subsection = {format = {\heiti \zihao{4}}},
  subsubsection = {format = {\heiti \zihao{-4}}},
}

% ---------- 代码样式（小字号等宽、带行号与边框） ----------
\lstset{
  language        = Python,
  basicstyle      = \ttfamily\zihao{6},
  keywordstyle    = \color{blue!70!black},
  commentstyle    = \color{green!40!black}\itshape,
  stringstyle     = \color{red!60!black},
  numbers         = left,
  numberstyle     = \tiny\color{gray},
  numbersep       = 6pt,
  frame           = single,
  framerule       = 0.4pt,
  rulecolor       = \color{gray!55},
  breaklines      = true,
  breakatwhitespace = false,
  columns         = fullflexible,
  keepspaces      = true,
  showstringspaces = false,
  tabsize         = 4,
  xleftmargin     = 1.8em,
  framexleftmargin = 1.8em,
  aboveskip       = 0.6em,
  belowskip       = 0.8em,
}

\begin{document}

\section*{附录：程序核心代码片段}

\noindent 核心步骤摘录；完整源程序见支撑材料（\texttt{q1--q4\_solver.py}、\texttt{q78\_sensitivity.py}、\texttt{verify\_q*.py}）。
全套参数：$R=0.02$ m，$h=25$ W/(m$^2\!\cdot$K)，$h_m=8\times10^{-7}$ m/s。

\subsection*{1\quad 求解内核}

\subsubsection*{1.1\quad 三对角追赶法}

\begin{lstlisting}
def thomas(a, b, c, f):
    n = len(f); p = np.zeros(n); q = np.zeros(n)
    p[0] = c[0] / b[0]; q[0] = f[0] / b[0]
    for i in range(1, n):
        m = b[i] - a[i] * p[i - 1]; p[i] = c[i] / m; q[i] = (f[i] - a[i] * q[i - 1]) / m
    x = np.zeros(n); x[-1] = q[-1]
    for i in range(n - 2, -1, -1): x[i] = q[i] - p[i] * x[i + 1]
    return x
\end{lstlisting}

\subsubsection*{1.2\quad 守恒型有限体积离散 + Crank--Nicolson 单步}

\noindent 随体坐标 $\zeta$ 下（$R$ 取常数即退化为径向格式），$C_v$ 为容量系数，
$\beta$ 为表面 Robin 系数（热 $hR/k_s$、质 $h_mR/D_s$）：

\begin{lstlisting}
def crank(u, K, Cv, beta, b0, b1, dt, R, N, dz):
    Km = (K[:-1] + K[1:]) / 2; r2 = R * R
    a = np.zeros(N); b = np.zeros(N); c = np.zeros(N)
    b[0] = -4 * Km[0] / (dz ** 2 * r2); c[0] = -b[0]
    for i in range(1, N - 1):
        a[i] = Km[i - 1] * (i - 0.5) / (i * dz ** 2 * r2)
        c[i] = Km[i] * (i + 0.5) / (i * dz ** 2 * r2)
        b[i] = -(a[i] + c[i])
    v = (1.0 - ((N - 1.5) * dz) ** 2) / 2.0; B = beta * K[-1]
    a[N - 1] = Km[N - 2] * (N - 1.5) / (v * r2)
    b[N - 1] = -(a[N - 1] + B / (v * r2)); c[N - 1] = 0.0
    s = np.zeros(N); s[N - 1] = B / (v * r2)
    L = lambda y: a * np.roll(y, 1) + b * y + c * np.roll(y, -1)
    f = Cv / dt * u + 0.5 * L(u) + 0.5 * (b0 + b1) * s
    return thomas(-0.5 * a, Cv / dt - 0.5 * b, -0.5 * c, f)
\end{lstlisting}

\subsubsection*{1.3\quad 交错 Picard 耦合步}

\begin{lstlisting}
def picard(T, C, t0, t1, dt, R, N, dz, rho, cp, k, D):
    Tn, Cn = T.copy(), C.copy()
    for it in range(1, 21):
        K = k(C)
        T = crank(Tn, K, rho(C) * cp(C), h * R / K[-1], Tair(t0), Tair(t1), dt, R, N, dz)
        Dc = D(C, T + 273.15)
        C = crank(Cn, Dc, np.ones(N), hm * R / Dc[-1], Cair(t0), Cair(t1), dt, R, N, dz)
        if it > 1 and max(np.max(np.abs(T - Tp)) / max(np.max(T), 1.0),
                          np.max(np.abs(C - Cp)) / max(np.max(C), 1.0)) < 1e-8:
            break
        Tp, Cp = T.copy(), C.copy()
    return T, C
\end{lstlisting}

\subsection*{2\quad 物性与离散参数}

\begin{lstlisting}
# 附录2（问题1：常数物性 + 非线性 D）
rho2, cp2, k2 = 820.0, 2600.0, 0.36
D2 = lambda C: 7e-9 * np.exp(-0.89 / C)

# 附录3（问题2、3）
rho3 = lambda C: 650.0 + 128.0 * C
cp3 = lambda C: 1450.0 + 2736.0 * C / (C + 1.0)
k3 = lambda C: 0.21 + 0.38 * C / (C + 1.0)
D3 = lambda C, T: 2.4e-3 * np.exp(-0.45 / C) * np.exp(-3850.0 / T)

# 附录4（问题4）
rho4 = lambda C: 760.0 + 90.0 * C
cp4 = lambda C: 1850.0 + 2150.0 * C / (C + 1.0)
k4 = lambda C: 0.12 + 0.20 * C / (C + 1.0)
D4 = lambda C, T: 4.2e-4 * np.exp(-0.30 / np.maximum(C, 0.02)) * np.exp(-3850.0 / T)
\end{lstlisting}

\begin{table}[htbp]
\centering
\zihao{6}
\setlength{\tabcolsep}{4pt}
\renewcommand{\arraystretch}{1.25}
\begin{tabular}{@{}lllll@{}}
\toprule
项目 & 问题 1 & 问题 2 & 问题 3 & 问题 4 \\
\midrule
物性 & 附录 2 & 附录 3 & 附录 3 & 附录 4 \\
空间离散 & $\Delta r=0.03125$ mm，$N=641$ & $\Delta r=0.125$ mm，$N=161$
         & $\Delta r=0.0625$ mm，$N=321$ & $\Delta\zeta=1/160$，$N=161$ \\
时间步长 & $\Delta t=0.125$ s & $\Delta t=0.25$ s & $\Delta t=2.5$ s & $\Delta t=5$ s \\
几何 & $R=2$ cm 固定 & 固定 & 固定 & $R(t)$ 附件 2 \\
求解区间 & 0--1800 s & 0--10800 s & 至判据 & 至判据 \\
\bottomrule
\end{tabular}
\end{table}

\subsection*{3\quad 烘房条件与收缩动边界}

\begin{lstlisting}
Tair = lambda t: np.interp(t, d1[:, 0], d1[:, 1]) if t <= 14400 else 50.165
Cair = lambda t: np.interp(t, d1[:, 0], d1[:, 2]) if t <= 14400 else 0.04986

_Rp = PchipInterpolator(d2[:, 0], d2[:, 1] / 100.0, extrapolate=False)   # 附件2 R(t)
RofT = lambda t: _Rp(t) if t <= 259200 else 0.01198
\end{lstlisting}

\subsection*{4\quad 烘干时长：判据与事件求根}

\begin{lstlisting}
CT = 0.15
for n in range(1, nmax + 1):
    Cp_ = C.copy()
    T, C = picard(T, C, (n - 1) * dt, n * dt, dt, R, N, dz, rho4, cp4, k4, D4)
    if C.max() <= CT:                      # g=max C-CT 在相邻时间层间线性求根
        g = (Cp_.max() - CT) / (Cp_.max() - C.max())
        tstar = (n - 1 + g) * dt
        Cstar = Cp_ + g * (C - Cp_)
        break
\end{lstlisting}

\subsection*{5\quad 固定距离采样（表 6）}

\begin{lstlisting}
def interp_dist(C, R, dj):
    return float(PchipInterpolator(np.arange(len(C)) * dz, C)(dj / R))
\end{lstlisting}

\subsection*{6\quad 验证}

\subsubsection*{6.1\quad Bessel 级数解析锚点（问题 1）}

\begin{lstlisting}
f = lambda b: b * j1(b) - h * R / k2 * j0(b)
bs = np.array([brentq(f, s, s + 0.5) for s in np.arange(0.05, 200, 0.5) if f(s) * f(s + 0.5) < 0][:40])
An = 2 * j1(bs) / (bs * (j0(bs) ** 2 + j1(bs) ** 2))
Texact = lambda x, t: Ta + (T0 - Ta) * (An * np.exp(-k2 / (rho2 * cp2) * bs ** 2 * t / R ** 2)
                                         * j0(bs * x / R)).sum()
\end{lstlisting}

\subsubsection*{6.2\quad 离散守恒恒等式}

\begin{lstlisting}
v = np.zeros(N); v[0] = dz ** 2 / 8
v[1:N - 1] = np.arange(1, N - 1) * dz ** 2
v[N - 1] = (1.0 - ((N - 1.5) * dz) ** 2) / 2.0
W = lambda C, R: R * R * (C * v).sum()          # 校验 W(C_f)-W(C_0) = -R∫h_m(C_R-C_air)dt
\end{lstlisting}

\subsubsection*{6.3\quad 烘干时长收敛链与 Richardson 外推}

\begin{lstlisting}
t = np.array([tstar(dz) for dz in (1 / 80, 1 / 160, 1 / 320)]) / 3600
p = np.log2((t[1] - t[0]) / (t[2] - t[1])); ext = t[2] + (t[2] - t[1]) / (2 ** p - 1)
\end{lstlisting}

\subsection*{7\quad 片段与源程序对照}

\begin{table}[htbp]
\centering
\zihao{-5}
\begin{tabular}{@{}ll@{}}
\toprule
片段 & 所在源程序 \\
\midrule
1.1--1.3 & 四个求解器共用 \\
2--5 & \texttt{q1}--\texttt{q4\_solver.py}、\texttt{q78\_sensitivity.py} \\
6.1--6.2 & \texttt{q1\_solver.py}、\texttt{verify\_q1.py}、\texttt{q3/q4\_solver.py} \\
6.3 & \texttt{verify\_q3.py}、\texttt{verify\_q4.py} \\
\bottomrule
\end{tabular}
\end{table}

\end{document}
